% TOP_PROBLEM: {"problemId": "problem.benjamini-schramm-conjecture-p-c-p-u-for-non-amenable-cayley-graphs", "canonicalTitle": "Benjamini-Schramm conjecture p_c<p_u for non-amenable Cayley graphs", "releaseRank": 360, "primaryDomain": "probability_ergodic_dynamics", "primaryDomainLabel": "Probability, ergodic theory & dynamics", "displayStatus": "open", "releaseStatus": "open"}
% !TeX program = lualatex
% Definition and sources only; this file is not an LLM solution attempt.
\documentclass[11pt]{article}
\usepackage[margin=25mm]{geometry}
\usepackage{fontspec}
\setmainfont{Latin Modern Roman}
\usepackage{amsmath,amssymb,mathtools}
\usepackage{unicode-math}
\setmathfont{Latin Modern Math}
\usepackage{xurl,hyperref}
\hypersetup{colorlinks=true,urlcolor=blue}
\setcounter{secnumdepth}{0}
\setlength{\parindent}{0pt}
\setlength{\parskip}{5pt}
\setlength{\emergencystretch}{3em}
\begin{document}
\section*{\texorpdfstring{Benjamini-Schramm conjecture p\_c<p\_u for non-amenable Cayley graphs}{Benjamini-Schramm conjecture p\_c<p\_u for non-amenable Cayley graphs}}\label{360}
\subsection{Definitions and mathematical statement}
For a locally finite Cayley graph $G$, Bernoulli bond percolation independently retains each edge with probability $p$. Let $N_\infty$ be the number of infinite open clusters and let $o$ be a fixed vertex. Define
\[
 p_c=\inf\{p:{\mathbb{P}}_p(o\leftrightarrow\infty)>0\},\qquad
 p_u=\inf\{p:{\mathbb{P}}_p(N_\infty=1)=1\}.
\]
A finitely generated group is amenable if it has a translation-invariant mean, equivalently a F\o lner sequence of finite sets; nonamenability is failure of this property. The selected assertion is $p_c(G)<p_u(G)$ exactly when the infinite underlying group is nonamenable, for every finite generating set defining a Cayley graph. The thresholds are infima, so the statement does not prescribe cluster behavior at either endpoint. Site percolation and arbitrary nontransitive graphs are not silently substituted.
\par\smallskip\textit{Formulation and concept references:} \href{https://ar5iv.labs.arxiv.org/html/1712.04651}{[S1]}.

\subsection{Short English statement}
Does a Cayley graph have a nonempty percolation phase with multiple infinite clusters exactly when its group is nonamenable?
\subsection{Sources}
\begin{itemize}
\item[S1]H. Duminil-Copin, "Sixty years of percolation", Proc. International Congress of Mathematicians 2018, vol. IV, pp. 639-672; arXiv:1712.04651.
\url{https://ar5iv.labs.arxiv.org/html/1712.04651}.
\textit{Formulation source.}
\end{itemize}
\end{document}
